\section*{Step 253: de Bruijn--Newman Carrier}

\paragraph{Carrier.}
Let
\[
H_t(z)=\int_0^\infty e^{t u^2}\Phi(u)\cos(zu)\,du,
\]
where
\[
\Phi(u)=\sum_{n\ge1}
\bigl(2\pi^2n^4 e^{9u}-3\pi n^2 e^{5u}\bigr)
e^{-\pi n^2e^{4u}}.
\]
This is the heat-flow deformation of the Riemann xi function on the critical
line.  The de Bruijn--Newman constant is
\[
\Lambda=\inf\{t\in\mathbb R: H_t \text{ has only real zeros}\}.
\]

\paragraph{Classical status.}
de Bruijn's 1950 theorem gives the heat-flow real-zero persistence direction:
once real-rootedness holds at some time, it holds for later times.  Newman
proved in 1976 that a finite threshold \(\Lambda\) exists and conjectured
\(\Lambda\ge0\).  Rodgers--Tao proved \(\Lambda\ge0\).  Polymath15 proved
the unconditional numerical upper bound \(\Lambda\le0.22\).

\paragraph{RH equivalence.}
The Riemann Hypothesis is equivalent to \(\Lambda\le0\).  Combining this with
Rodgers--Tao's lower bound, the remaining target is exactly
\[
\Lambda=0.
\]

\paragraph{CRCFT mode.}
The carrier is CRE and classifies as CRCFT--TE: closing the residual is
target-equivalent to RH.  It is structurally distinct from the Burnol--Sonine
operator carrier, but no bridge is needed for this classification because the
de Bruijn--Newman equivalence is direct.

\paragraph{Verdict.}
\[
\mathrm{V\_de\_Bruijn\_Newman\_CRCFT\_TE}.
\]
No proof of RH is claimed.
